\section{Introduction}
\label{sec:intro}

Text detection on under-resourced scripts is a persistent bottleneck for
real-world document understanding. Mongolian, written vertically and used
across both handwritten archival pages and modern signage, is a
representative case: a single script appears in two visually disjoint
domains, and supervised training data exists in different quantities for
each. Practitioners therefore face an asymmetric supervision problem in
which a model trained on the dominant domain must transfer to a target
domain that may share script but differs sharply in background,
illumination, layout, and writing style. Naive cross-domain transfer
of a single-domain detector typically yields target-domain
$\hmean$ that is one to two orders of magnitude lower than the
in-domain oracle. The gap is severe enough that in one of the two
transfer directions we study below, the source-only model produces
zero true positives at every post-process threshold setting.

Unsupervised domain adaptation (UDA) offers an attractive remedy
because it requires no target-domain labels and can be bolted on top of
an existing detection backbone. The canonical UDA recipe for
object/text detection combines two ingredients: feature-level
distribution alignment via a gradient-reversal domain classifier
(DANN~\cite{ganin2015dann}) and target-domain self-training via an
EMA teacher generating pseudo-labels for the student~\cite{li2022adapttch}.
Each ingredient is light, simple to implement, and adds only a small
runtime overhead, making the combination especially appealing for
small-corpus, low-budget settings such as ours (Apple-Silicon MPS,
single device, $\sim$2500 total images, ResNet-18 backbone).

Despite the simplicity of the recipe, applying it to a real cross-domain
benchmark surfaces failure modes that the literature has largely
glossed over. We report three. First, when the source-only model
produces effectively zero detections on the target domain, the
randomly initialised teacher generates pseudo-labels indistinguishable
from random predictions, and the self-training loop never escapes the
trivial fixed point of predicting nothing -- a phenomenon we call
\emph{cold-start collapse}. Second, on the same UDA pipeline applied
to both transfer directions, the magnitude of improvement varies by an
order of magnitude: the harder direction (scene$\to$handwriting,
source-only $\hmean = 0$) gains $+6.26$~pp $\hmean$ on the target test
split, while the easier direction
(handwriting$\to$scene, source-only $\hmean \approx 0.075$ after
threshold tuning) gains only $+0.14$~pp -- well within run-to-run
variance. We call this the \emph{asymmetric UDA gain phenomenon}.
Third, contrary to the usual assumption that more epochs help, both
directions exhibit a sharp \emph{5-epoch structural sweet spot}:
extending training to 25 epochs causes teacher--student drift
collapse and degrades target $\hmean$ by an order of magnitude even
under a conservative pseudo-label weight schedule.

These observations suggest that practitioners of self-training UDA on
under-resourced scripts should re-examine three default assumptions:
that random initialisation suffices, that UDA produces a roughly
uniform gain across transfer directions, and that more training
epochs are safer than fewer. The remainder of this paper presents the
pipeline, diagnoses each failure mode, proposes a single-line fix for
the cold-start case, and quantifies all three findings on a clean
Mongolian cross-domain benchmark.

\paragraph{Contributions.}
\begin{enumerate}
  \item We diagnose a \emph{cold-start collapse} failure mode in
        teacher--student self-training UDA when the source-only model
        yields near-zero target-domain detections, and propose a
        single-line remedy: initialise both student and teacher from the
        source-only checkpoint.
  \item We empirically establish an \emph{asymmetric UDA gain}
        phenomenon on a Mongolian cross-domain benchmark: UDA produces
        a $+6.26$~pp $\hmean$ jump on the harder direction
        (scene-to-handwriting, where source-only is structurally
        broken) and only $+0.14$~pp on the easier one
        (handwriting-to-scene, where source-only already has a weak
        target-domain signal).
  \item We document a \emph{5-epoch structural sweet spot}: extending
        training to 25 epochs causes teacher--student drift collapse in
        \emph{both} directions, even under a conservative pseudo-label
        weight ramp.
  \item We release a clean, MPS-compatible reproducible pipeline
        (DBNet ResNet-18 + EMA teacher + DANN) with all
        train/eval/sweep/visualisation tools and 5 trained checkpoints.
\end{enumerate}
